\chapter{George Dantzig}

George Dantzig’s development of the Simplex method for linear programming models represents a watershed moment in the history of industrial computing that led directly to major commercial successes for the digital computer and altered the course of the oil industry in Texas and its academic partners at UT Austin and Rice University in Houston. The commercialization of Dantzig’s Simplex functioned as a primary economic engine for the fledgling digital computer industry, specifically by solving the blending problem inherent in petroleum refining. By the mid-fifties, the Texas oil industry recognized that linear programming models could be used to optimize the mixture of crude oil and additives to produce specific grades of gasoline at the lowest possible cost. Before the fifties, petroleum refineries relied on manual calculations and trial-and-error to determine how to mix crude oil fractions, like heavy naphtha and catalytic reformate, to create gasoline. This manual approach frequently resulted in "quality giveaway", a wasteful practice where companies over-blended expensive, high-octane components to ensure the final fuel met minimum quality standards. Because Dantzig's linear programming models could identify the absolute best mathematical outcome rather than just a good one, it eliminated this waste and fundamentally redefined refinery efficiency.

The transition from manual calculations to algorithmic optimization compelled companies like Texaco, Humble Oil (Exxon), and Shell to become early major private-sector purchasers of mainframe computers, effectively subsidizing the research and development costs of the entire computing field. The Texas oil industry's need to solve both linear programming models and finite difference subsurface models drove the commercial market for the earliest mainframes, such as the IBM 701 in 1952 and the IBM 704 in 1954. By paying the astronomical rental fees for these machines, often exceeding \$15,000 per month, these companies provided the vital capital IBM needed to transition from experimental prototypes to reliable, mass-produced computers. 

This surge in industrial investment fundamentally reshaped the academic landscape of Texas, creating a robust research corridor between Houston and Austin. UT Austin and Rice became centers of excellence for numerical analysis, as oil companies funded the establishment of computation centers to train the specialized workforce required to manage these new digital assets. The computational centers and mainframes used for refinery and subsurface modeling were perfectly suited for the complex trajectory and dynamics calculations required for spaceflight. For instance, the IBM 704 was utilized to track satellite orbits in 1957, and this pre-existing computational excellence at Rice University supported NASA's decision to move the Manned Spacecraft Center to Houston in the early sixties. This era established the Texas model of high-tech development as a feedback loop where industrial profit from the energy sector was reinvested into university-based computational research, providing the intellectual infrastructure that would later support the state's aerospace and semiconductor booms. 

During World War II, operating within the Army Air Forces Statistical Control Division, an elite unit that included several of the future so-called Whiz Kids, Dantzig was tasked with quantifying the aerial war effort. In many ways the Statistical Control Division was the wartime precursor to the RAND corporation, where Dantzig would work in the fifties. He held his senior planning role as a civilian, managing data collection on sorties flown, bombs dropped, and aircraft losses. He assisted in creating “Programs”, highly detailed logistical plans for procuring airplanes, manufacturing parts, and training personnel. His exceptional service during this period was defined by his ability to synchronize the chaotic variables of global combat into a coherent logistical framework, a feat that necessitated the invention of the very planning structures that he would later automate with the Simplex method. These highly detailed logistical plans required the simultaneous coordination of industrial production, personnel training, and frontline deployment, creating a system of systems that pushed the limits of human calculation. 

Dantzig recalled that programming the war effort was essentially "doing the same thing as planning a whole country," requiring the coordination of hundreds of thousands of material goods and up to fifty thousand personnel specialties, and he became a skilled expert at managing these massive operations using only hand-operated desk calculators and primitive punched-card equipment. It was the challenge of managing these massive, interlocking schedules by hand that convinced him of the necessity for a universal mathematical method capable of modeling such systems. One of his key insights was why manual planning was so chaotic and mathematically limited. Before linear programming, planners had no way to computationally determine the best combination of actions for an overall objective. Instead, they relied on hundreds of highly subjective ground rules dictated by commanders. For example, a rule might simply be to "Go ask General Arnold which alternative he prefers," forcing planners to constantly spend time seeking resolutions rather than mathematically optimizing a clear goal. 

Dantzig later identified one of his greatest contributions as the recognition that these convoluted ground rules could be replaced by a general objective function. His transition from managing manual systems to inventing linear programming was directly triggered by a post-war retention effort. In 1946, Dantzig was preparing to leave the Pentagon for an academic position at UC Berkeley. To keep him, his Air Force colleagues offered him a much higher salary and challenged him to mechanize the planning process to allow for rapid, last-minute recalculations of time-staged force deployments and supply logistics. This challenge to automate what was previously done by hand culminated in his 1947 formulation of linear programming models and the Simplex Method.

In the late forties and early fifties, alongside the birth of the computer, there was a tectonic shift in industrial management, transitioning from the ground rules of adhoc manual planning to the rigorous mechanization of planning programs. The era was marked by escalating tensions that created an unprecedented demand for advanced computational methods. The sheer scale of global mobilization, coordinating hundreds of thousands of material types, transportation routes, and personnel specialties, exposed the fragility of intuitive, rule-of-thumb logistics. World War II exposed the fatal flaws of traditional, decentralized administrative bureaus that relied on intuition and manual ledger entries. To coordinate an air arm of over 2.4 million personnel and 80,000 aircraft, the War Department established Statistical Control in 1942 to impose a rationalized, top-down order. By decomposing complex military processes into quantifiable units, this group based supply requirements on computed planning factors, like the ratio of gasoline consumption to combat sorties. This rigorous quantification yielded massive efficiencies, saving billions of dollars in planned production during the war, and serving as the essential catalyst for the modern era of industrial efficiency. 

In 1946, ten Statistical Control veterans including future Secretary of Defense Robert McNamara and future Litton Industries founder Charles Thornton, offered their services as a package deal to Henry Ford II. At the time, Ford was hemorrhaging money under chaotic management; the company lacked basic cost accounting and actually estimated its expenses by weighing sacks of invoices. The group, soon dubbed the Whiz Kids, radically transformed Ford by introducing a management-by-numbers philosophy. They built a system of rigorous financial controls, implemented the company's first cash forecasts and capital budgets, and restructured the corporation so that each division functioned as an accountable profit center.

The technocratic paradigm pioneered by the Whiz Kids rapidly migrated to other complex industries, most notably the petroleum sector. Companies like Exxon developed corporate cultures deeply obsessed with mathematics, engineering, and systems optimization. This approach redefined the modern corporation as a "system of systems" with complex, networked, and distributed operations managed via mathematical models. At Exxon, this model-based resource allocation extended beyond physical supply chains and gasoline blending into sophisticated financial engineering and global risk management.
The institutionalization of mechanized planning fostered a business culture that equated mathematical modeling with operational reality. However, the historical record also highlights the dangers of this rigid technocratic paradigm. The strict reliance on quantifiable data created a monoculture that frequently ignored qualitative human realities. At Ford, for example, the Whiz Kids' strict doctrine of rejecting any investment that could not prove immediate bottom-line profit led to a severe degradation in product quality and innovation, as they refused to invest in new manufacturing equipment. Their over-systematization risked creating a mindset where unquantified factors simply ceased to exist for decision-makers, demonstrating that while model-based resource allocation is immensely powerful, it possesses blind spots when intuition and qualitative value are entirely removed from the equation. The mechanization of planning did more than simply improve industrial efficiency. It established a technocratic paradigm that would define the modern corporation and permanently tether industrial success to the efficacy of algorithmic resource allocation. 
This shift in management philosophy was perfectly timed with the arrival of digital computing. The transition from manual calculations to stored-program computers, like the UNIVAC I installed at the Pentagon in 1952, effectively automated the role of logistical models. Computers were no longer just calculators but became active participants in designing organizational strategy, allowing decision-makers to synchronize disparate production cycles with unprecedented precision. 

Dantzig’s work was foundational to the new fields of Operations Research, Systems Analysis, and Systems Engineering. Originating in England, where it was referred to as Operational Research, OR involved using mathematics to solve specific tactical problems during the thirties and forties, initially to integrate new early-warning radar technology into air defense. These early OR teams were radically interdisciplinary, employing physiologists, astrophysicists, and mathematicians to analyze human-machine interactions under combat conditions. They successfully applied statistical analysis to optimize convoy densities against U-boats and bomber formations to minimize losses. However, these British applications were focused strictly on the optimal use of existing equipment for specific, short-term missions.
Before Dantzig’s formulation of the Simplex Method, the management of large-scale organizational problems was largely an exercise in descriptive heuristics. By introducing a definitive algorithmic approach to linear programming, Dantzig enabled the quantification of optimal outcomes within complex, multi-variable environments. This breakthrough allowed the RAND Corporation and the Department of Defense to move beyond intuitive planning toward a model of Systems Analysis, where the rigorous comparison of competing strategic alternatives could be conducted with mathematical certainty. Instead of asking "How do we bomb this target efficiently?", Systems Analysis asked "Is bombing this target the most cost-effective way to win the war compared to other options?" It weighed costs against benefits across an entire system. While traditional OR asked how to execute a mission efficiently, Systems Analysis was future-oriented. It sought to find the optimal mix of forces, strategies, and unbuilt equipment to maximize national security subject to strict budget constraints. This cost-benefit approach eventually culminated in the Planning, Programming, and Budgeting System instituted by Robert McNamara, officially shifting the Pentagon toward a model of centralized, mathematically-informed decision making. 

The principles of modeling and optimization pioneered by Dantzig were essential to the formalization of Systems Engineering as a distinct technical discipline. The Air Force's Intercontinental Ballistic Missile program, which engaged more than 2,000 industrial contractors and over 40,000 personnel, was considered vastly more complex than even the Manhattan Project. To manage this, the Air Force hired the Ramo-Wooldridge Corporation to provide "systems engineering and technical direction". Systems engineers utilized mathematical optimization to perform critical tradeoff studies across the entire ICBM lifecycle. For example, if a propulsion system fell short of its thrust requirements, optimization models determined exactly how weight could be safely reduced in the guidance or re-entry subsystems to compensate. This ensured that the missile functioned as a harmonized whole despite the chaotic variables of simultaneous, concurrent development. The mathematical frameworks established by Dantzig allowed systems engineers to model the entire lifecycle of a project as a unified set of constraints and objectives, ensuring that individual technical successes translated into a functional whole. Ultimately, the Dantzig paradigm transitioned from the military-industrial complex to the commercial sector, where it continues to power the model-based algorithmic dispatch, facility location, and route optimization of modern logistics giants like Amazon and FedEx. 

\section{Simplex Linear Programming}

In 1947, while serving as a mathematical advisor to the U.S. Air Force Comptroller, Dantzig formulated the linear programming model as a system of linear inequalities, and created the Simplex Method of automatically solving the resulting linear system. Linear programming was one of the earliest major applications of automated linear algebra and was born just before, and in conjunction with, the first stored-program computers. Unlike many contemporary mathematical developments, LP was uniquely machine-ready, emerging in a symbiotic relationship with the first generation of computers. This parallel growth ensured that as the physical architecture of the computer evolved, it was immediately directed toward the most complex logistical and strategic challenges of the early Cold War. 

The Standards Eastern Automatic Computer, built by the National Bureau of Standards, was one of the first fully operational electronic stored-program computers in the US. Dantzig's work within the Air Force's Project SCOOP directly financed and catalyzed its development, with the Air Force providing \$400,000 to the National Bureau of Standards to develop the SEAC. Dantzig was involved with the practical aspects of using it to solve the first computerized LP problem. In 1951, under the direction of Alex Orden and NBS staff, the SEAC successfully solved the first electronic LP problem, a 48-equation, 71-variable Air Force deployment model, in an 18-hour run. In other words, a commercially significant, operational real-world LP problem was being run in 1951, at a time when other stored-program computers were still in development. Dantzig and the National Bureau of Standards demonstrated the viability of computerized linear algebra for real-world applications at a time when the field was still largely experimental. These events established a precedent for the projects of the following decade, proving that the digital computer could function as an indispensable tool for national and industrial decision making. The 1951 SEAC run effectively inaugurated a new era of high-performance computing, cementing the role of the mathematician-programmer as a central architect of the post-war technological order.
This milestone was rapidly followed by the Air Force installing the second UNIVAC I in the Pentagon basement in 1952, which handled deployment models containing up to 250 equations and 500 variables. To process these complex algorithms, early computer scientists utilized highly innovative software and memory management techniques. For instance, William Orchard-Hays developed a "collapsed vector" method for the IBM 701 to squeeze the data of sparse LP matrices into the machine's tiny 2,048-word memory, maximizing its computational usefulness. These early workarounds and hacks validated the potential for larger-scale numerical models, laying the groundwork for the massive aerospace and ICBM systems engineering projects later in the fifties. 

In the 1940s, the military used the word "program" to describe a proposed schedule of actions, training plans, logistical supply chains, or force deployments. Dantzig and his contemporaries distinguished this from coding, which referred to the actual instructions written for a machine. Linear programming, therefore, was the mathematical structure of the plan itself, a system of linear inequalities paired with a linear objective function to be optimized. LP provided a mathematical framework for optimizing a given objective function subject to a set of linear constraints. This allowed decision-makers to pursue an explicit goal, such as minimizing costs or maximizing combat sorties, subject to available resources and equipment capacities.

The Simplex method's revolutionary impact lay in providing, for the first time, a computationally feasible method for finding the single best solution among a potentially infinite number of possibilities. Before Simplex, finding the single optimal solution for operational-scale problems was mathematically impossible due to the astronomical number of potential combinations. Dantzig often illustrated this by noting that solving a simple assignment of 70 men to 70 jobs by checking every permutation by brute force would require "a solar system full of nano-second electronic computers running from the time of the big bang until the time the universe grows cold". Dantzig's algorithm revolutionized this by utilizing the geometric properties of the feasible region, which forms a multidimensional shape known as a convex polytope. Because mathematical theory proved that the optimal solution must lie at one of the vertices or corner points of this shape, Simplex provided a systematic, iterative procedure to move along the edges from one vertex to an adjacent one, continuously improving the objective value at each step until the absolute best outcome was identified.
This breakthrough had an immediate and profound impact on the American military-industrial complex, completely replacing subjective, ad hoc ground rules with algorithmic certainty. By modeling and automating planning, Dantzig's work helped establish a new cultural and technocratic paradigm where mathematical modeling was directly equated with operational reality. This rigorous, quantifiable approach to resource allocation spread rapidly from the Pentagon to major corporations like Ford and Exxon, redefining modern management as a system of systems and inextricably linking both industrial and military success to the power of numerical modeling and optimization. 

Up until this time, and during the birth of LP, the word computers often referred to rooms filled with desktop mechanical calculators operated manually, often by women. This is the reason that with the coming of the electronic stored-program computer, the first generation of computer programmers were often women. As the field transitioned to stored-program computers, women like Leola Cutler and Ida Rhodes became highly influential, acting as the "grandmothers of a great many programmers" and pioneering strict coding discipline. The task of coding, initially perceived as an extension of manual calculation, was often assigned to traditional human computers, but the complexity of implementing Simplex necessitated a disruption of this pattern. Dantzig emerged as a pioneering figure in this regard as well, who purposefully blurred the boundaries between human computers, coders, and modelers, recognizing that the successful execution of LP required a rigorous, hands-on synthesis of high-level mathematical modeling and low-level machine instructions.  

He was one of the early examples of a researcher who purposefully combined applied mathematics with hands-on coding. This was illustrated by the fact that he took LP and the Simplex method with him, moved to RAND in Los Angeles, and personally helped implement his code on the first generations of IBM computers from 1952 with the help of a few like-minded hands-on researchers. William Orchard-Hays was the primary coder who translated Simplex into a tight and efficient system at the RAND. Orchard-Hays brought the imagination and rigorous technical creativity required to implement complex mathematics on the primitive first generations of IBM hardware. Because this work preceded high-level languages like Fortran, Orchard-Hays and his team worked with very crude assemblers, writing their code in octal. By personally directing the coding of Simplex, Dantzig ensured that the mathematical nuances were preserved during the translation to binary logic. This work at RAND effectively transformed the digital computer into a dynamic laboratory for modeling, establishing the technical and professional standards that would later define the high-performance computing ecosystem in both California and Texas. 

In the early fifties, the Los Angeles environment represented a unique convergence of geopolitical urgency, industrial concentration, and the growth of a technological culture that would soon spread and proliferate. It was characterized by a dense network of interconnected entities, all oriented towards the practical real-world challenges of the era. The ecosystem was not just a collection of companies, factories, and universities, but a sophisticated knowledge network where the boundaries between academic inquiry, industrial production, and military command were deliberately porous. In 1954, the Air Force established the Western Development Division in Inglewood under Bernard Schriever, with responsibility over the accelerated ICBM effort. Because traditional airframe manufacturers lacked the comprehensive scientific expertise required to integrate the radical new technologies of propulsion, guidance, and re-entry, Schriever took the unprecedented step of contracting with a private company, the newly formed Ramo-Wooldridge Corporation, later TRW and also based in Inglewood, to serve as the system's technical arm. Essentially, what would have been handled by the government in that past was outsourced to industry, and Ramo-Wooldridge was made responsible for Systems Engineering and Technical Direction over all other contractors, effectively bridging the gap between abstract physics and industrial hardware. 
This ecosystem was fueled by an unprecedented influx of federal defense capital, which compelled private contractors and research institutions to adopt a systems-oriented worldview. For the engineers and mathematicians operating within this environment, the challenge of the era was not merely the construction of hardware, but the development of a comprehensive methodology for managing the radical complexity of modern weaponry and space exploration. Together, the WDD and Ramo-Wooldridge pioneered the concurrency concept, a radical acquisition strategy where the design, testing, manufacturing, and even the construction of launch bases and training of crews occurred simultaneously rather than sequentially. This approach was heavily reliant on new management tools like the Program Evaluation and Review Technique and advanced computer modeling to track thousands of design changes across millions of components.

\section{RAND}

Central to this intellectual infrastructure was the RAND Corporation, established near the Santa Monica pier to serve as a specialized bridge between the Air Force and the scientific community. Originally launched in 1946 as Project RAND within the Douglas Aircraft Company in Santa Monica, it spun off into an independent nonprofit in 1948 to retain the civilian brainpower that had helped win World War II. This was the first of the breakthrough R\&D labs dedicated to modeling and simulation, as was only fitting for an organization that was named directly from the phrase Research and Development. The only comparable organizations at the time were Los Alamos, which was temporarily scaling down in the immediate post war years, and perhaps Bell Labs. The 1952 RAND milieu institutionalized the concept of R\&D as a unified, experimental process rather than a linear sequence of events. 

RAND was the intellectual bridge within the aerospace ecosystem. Rather than simply asking how to achieve a target efficiently, RAND analysts used numerical models, linear programming, and game theory to ask whether achieving that target was the most cost-effective way to balance national security goals against alternative strategies. Meanwhile, RAND scientists from institutions like Caltech and MIT were deeply integrated into the military-industrial efforts, serving on critical advisory panels like the Teapot Committee that completely redefined the Atlas ICBM's technical specifications. The tactical, mission-specific Operational Research of World War II was transformed into the much broader, future-oriented discipline of Systems Analysis. Its primary mission, the systematic incorporation of technical analysis into military strategy, precipitated a revolution in the social and physical sciences. By institutionalizing the use of quantitative modeling and operations research, RAND provided the strategic logic that justified the engineering efforts of the neighboring aerospace firms, TRW in particular. This environment created a specific breed of researcher-practitioner who was equally comfortable with abstract mathematical theory and the pragmatic constraints of hardware implementation, and it forged the analytical tools and professional standards that would define the trajectory of scientific computing and systems engineering.

To process massive logistical and strategic models, RAND became a global center for computer science, building the JOHNNIAC mainframe in its basement and developing new programming languages like SIMSCRIPT. The RAND of 1952 was a very early example, second only to Los Alamos and Project Whirlwind perhaps, of an environment with a sizable group of what were in effect some of the first computer hackers. While the organization was originally conceived as a strategic bridge between the Air Force and the scientific community, it quickly evolved into an environment where the distinction between theoretical research and hands-on implementation was aggressively dismantled. An example of this was the urgent need to operationalize George Dantzig’s linear programming models on the first available commercial hardware, the IBM CPC, 701, and 704. By engaging in the laborious task of machine-level coding, RAND researchers moved further beyond the era of human computers on the path that had already been blazed at Los Alamos, establishing a new professional identity centered on the creative and iterative manipulation of digital logic to solve real-world problems. 

The proximity of elite mathematicians like Dantzig to the actual hardware allowed for a rapid feedback loop between modeling theory and computational performance. This synergy transformed the digital computer from a static calculating device into a dynamic laboratory for modeling and simulation. The methodologies pioneered during this era, ranging from the development of custom diagnostic tools to the optimization of code for memory-constrained environments, laid the intellectual and cultural foundation for the high-performance computing centers that would later emerge, and the RAND of 1952 became a prototype of the hacker labs to come. 

RAND emerged as the preeminent intellectual architect of Cold War strategy and a multidisciplinary approach that fused rigorous quantitative methodology with national security policy via foundational developments in systems analysis and game theory. These were radical departures in strategic thinking that treated military planning not as an art but as a complex system of interconnected variables that could be computationally modeled and optimized, much like Dantzig’s LP problems. RAND became a hub for mathematics and game theory, enlisting John von Neumann and Oskar Morgenstern, alongside prominent researchers who pioneered dynamic programming to handle multistage decision processes. A 1954 RAND study on the selection and use of strategic air bases fundamentally altered U.S. nuclear doctrine. By analyzing base vulnerability and logistical supply chains, the study shifted the U.S. from a first-strike posture to a secure second-strike deterrence strategy, ensuring bombers could survive a surprise Soviet attack to retaliate, and ultimately saving the Air Force billions of dollars. This was all contemporaneous with Project Whirlwind’s conversion into the SAGE continental air defense system and Lincoln Labs, and the birth of Livermore Labs in competition with Los Alamos. In some sense Los Alamos, RAND, TRW, Lincoln, and Livermore all had much more in common than not, with perhaps the biggest distinction being that RAND and TRW were at least nominally in the private sector.

These methodologies became synonymous with RAND and were applied to a wide array of national security challenges. Rather than viewing defense planning as a qualitative art practiced by experienced commanders, RAND’s theorists modeled it as a set of variables that could be computationally optimized, establishing a logic of cost-effectiveness that would dominate federal procurement and strategic doctrine for decades. The culmination of this era was the publication of Charles Hitch and Roland McKean's The Economics of Defense in the Nuclear Age in 1960. This work argued that the economic use of scarce resources must be a critical aspect of defense planning. When Robert McNamara became Secretary of Defense in 1961, he appointed Hitch as Comptroller to implement the Planning, Programming, and Budgeting System. This officially cemented RAND's systems analysis and cost-effectiveness logic as the mandatory framework for building U.S. defense budgets and managing strategic planning. These methodologies required the systematic evaluation of vast decision spaces, where the consequences of various strategic choices, ranging from nuclear deterrence postures to logistical supply chains, were assessed through the lens of modeling and optimization. 

In 1954 RAND completed the construction of its own mainframe computer, the JOHNNIAC, a 40-bit word machine based on the von Neumann architecture under development at the Institute for Advanced Study. The JOHNNIAC was one of six related stored-program machines, the first and second being the ENIAC at IAS and MANIAC at Los Alamos. RAND installed Telemeter Magnetics core memory around late 1954 or early 1955, making JOHNNIAC the first operational computer with core memory in the world. This upgrade was a game-changer for reliability, making the machine at least five times more reliable than the notoriously temperamental IBM 701. As a von Neumann architecture machine, JOHNNIAC eschewed floating-point hardware. Rather than the floating-point arithmetic required for aerospace engineering, for example, the JOHNNIAC functioned as a general-purpose testbed for a series of software innovations that helped define the early computing landscape. It served as the primary incubator for the first heuristic programs and the nascent field of Artificial Intelligence, facilitating the transition from numerical calculation to automated logical reasoning. 

The RAND group utilized the machine to pioneer early timesharing systems, allowing multiple researchers to interact with the computer simultaneously, a radical departure from the batch processing model of the era. The JOHNNIAC was less a production machine than an experimental platform where the technical frameworks for artificial intelligence, advanced programming languages, and networked computing were first institutionalized. It was a general-purpose testbed for almost every major software innovation RAND produced for the next 13 years. Around 1960 to 1962, RAND engineers added typewriters and buffers to the machine to build the JOHNNIAC Open Shop System. This represented the world's first on-line timeshared computer system, introducing one of the earliest interactive programming languages for individual users. The JOHNNIAC continued to serve as a vital experimental platform until it was finally decommissioned in 1966, and many of its features were carried over in parallel to RAND’s spinoff Software Development Corporation and its SAGE AN/FSQ-32 where they directly influenced events at UT Austin, as will become apparent. 

RAND simultaneously advanced the technological frontiers via satellite reconnaissance in Project Feed Back, the pioneering of heuristic computing and artificial intelligence, and the application of psycho-political modeling to predict Soviet leadership dynamics. By the early sixties, RAND was arguably the world’s most advanced computer science laboratory. Because their strategic problems were too complex for existing commercial machines, they were forced to build their own hardware, invent their own programming languages, and effectively found the field of Artificial Intelligence. While the term Artificial Intelligence was coined in 1956, the actual work began at RAND a year prior with the collaboration of Allen Newell, Herbert Simon, and Cliff Shaw. Their research focused on "symbolic (non-numerical) processes to simulate human thinking". 

To manage the complex, branching logic of human problem-solving, Cliff Shaw pioneered the use of linked lists, a revolutionary data structure that permitted dynamic memory allocation. This led to the 1957 development of Information Processing Languages, the direct precursor to Lisp. These developments drove the first successful AI programs, and allowed for the creation of the Logic Theorist and General Problem Solver, which utilized means-ends analysis to decompose expansive objectives into a hierarchy of manageable sub-goals. The results included successful software to play chess, simulate ground combat, and perform early language-processing and Russian translation. By the time the first scholarly book on AI was published in 1963, six of its twenty chapters were previous RAND research reports. This pivot from calculation to cognition marked the birth of symbolic AI and the belief that the structural laws of human thought could be replicated through the algorithmic manipulation of non-numeric data. These innovations provided a framework for early AI research and began to reconfigure the American computer science curriculum. The methodologies established on the JOHNNIAC at RAND, specifically the use of recursive list processing, provided the direct intellectual lineage for the creation of Lisp, ensuring that the RAND model of symbolic processing would remain the standard for artificial intelligence research. Again, intensive work on Lisp continued at RAND’s spinoff SDC and the SAGE Q-32, work that eventually moved directly on to the UT Austin CDC 6600. 

In the early sixties, RAND researcher Paul Baran was tasked with the problem of maintaining control of the US nuclear arsenal. Baran’s research was motivated by the vulnerability of the nation’s centralized telecommunications infrastructure, which functioned as a fragile spider web susceptible to a single, well-placed nuclear strike. To mitigate this risk, Baran proposed a distributed network model that eliminated the necessity of primary switching centers. His radical solution, detailed in his 1964 report "On Distributed Communications," proposed breaking messages into discrete pieces. Instead of relying on a centralized hub, these pieces were sent into a network of connected points and allowed to "find their own way through" by bouncing to the next open node if any of the network was destroyed. By treating the network as a self-healing mesh of autonomous nodes, Baran transitioned communication theory from the rigid logic of the circuit to the fluid, resilient logic of the packet. This shift ensured that strategic command of the nuclear triad could be maintained through a landscape of partial systemic failure. This distributed architecture provided the foundational blueprint for the subsequent development of the Arpanet and global internet. 

Baran’s invention of discrete message blocks allowed for adaptive routing, where the network itself functioned as a decentralized intelligence capable of determining the optimal path for data in real-time. While his proposals were initially met with skepticism by the established telecommunications monopolies of the era, the technical viability of packet switching was validated by Donald Davies’s concurrent work at the National Physical Laboratory in the United Kingdom. Together, these innovations dismantled the traditional constraints of distance and centralized control, institutionalizing a new paradigm of digital communication that prioritized survivability and open-ended scalability. Despite authoring the foundational blueprint for one of the most transformative technologies in human history, Baran remained modest about his individual contribution. He famously compared the creation of the internet to the building of a grand edifice, noting that it was built by many hands over many years: “New people come along and each lays down a block on top of the old foundations, each saying, I built a cathedral”. 

RAND was also an early pioneer in human-computer interaction, and emphasized that computers should be accessible to researchers, not just specialists. The RAND Tablet or Grafacon, short for Graphic Converter, was marketed commercially by Bolt, Beranek and Newman and pushed these concepts even further. Developed in 1964, the tablet utilized a sophisticated grid of sensing wires to translate the analog motion of a stylus into digital coordinates with unprecedented accuracy. It was one of the first devices that permitted the direct input of handwritten text and freehand drawings into a computer. Demonstrating its unprecedented capabilities in digital handwriting recognition, researchers utilized the tablet to successfully look up hand-drawn Chinese characters. This device served as the primary ancestor to the modern digitizing tablet, enabling the first rigorous explorations of digital handwriting recognition and interactive mapping. By allowing users to input diagrams and hand-drawn data directly into the computer, the RAND Tablet provided the empirical foundation for the development of graphical user interfaces. 

Central to this transformation was the development of JOSS, the Johnniac Open Shop System , one of the world's first successful timesharing operating systems. Rather than using remote teletypes, this timesharing capability was achieved by adding typewriters and buffers to the JOHNNIAC mainframe around 1961 or 1962. By allowing multiple researchers to interact with the JOHNNIAC simultaneously, JOSS effectively democratized access to computational power. The system was famously characterized by its conversational interface, which utilized natural language error messages to guide non-specialist users through complex calculations. This emphasis on a user-friendly open shop environment dismantled the technocratic barrier between the scientist and the machine, institutionalizing the concept of the computer as a personal intellectual assistant rather than a centralized clerical tool. This work as well moved onward to RAND’s spinoff SDC and the SAGE Q-32, where it was encountered by Forest Baskett who would eventually carry it onward to the UT Austin CDC 6600.

\section{IBM CPC At RAND}

Before all of the other computing advances at RAND came Dantzig and Simplex in 1952. This was the heroic age, when the intellectual heft of Simplex collided with the extreme limitations of the early computers, and the IBM Card-Programmed Electronic Calculator was the battleground for this struggle. Unlike nascent stored-program computers like the ENIAC and JOHNNIAC, the CPC was a distributed assembly of punched-card machines that required constant human intervention to maintain the flow of logic. Dantzig and his key collaborator, William Orchard-Hays, didn't just implement and run Simplex. They were forced to reinvent it on-the-fly to fit the machine, and in the process to discover the pragmatic and ingenious hacks that became standard procedures for linear programming models. This was the ultimate in “code bumming” or the process of heavily optimizing a program or routine to make it as fast or compact as possible, often at the expense of code readability.

The IBM CPC was not much of a computer in the modern sense, and its development actually originated from clandestine, user-driven engineering at Northrop Aircraft rather than inside IBM's own research labs. In late 1946, a specialized computing group at Northrop led by engineers Greg Toben, Bill Woodbury, and Rex Rice was tackling complex aerospace calculations such as jet propulsion and guided missile trajectories, which vastly outstripped the capacity of standard accounting machines. Because true stored-program computers were not yet available, the Northrop team decided to build their own hardware solver by merging two leased IBM machines: the new IBM 603 Electronic Multiplier, which was fast but lacked sequencing control, and the older IBM 405 Accounting Machine, which was slow at math but had excellent card-reading and printing capabilities. In direct violation of their IBM rental agreements, the Northrop engineers took the protective covers off the machines, exposed their internal wiring, and physically linked the two units together. They affectionately dubbed their makeshift, hybrid creation the poor man's ENIAC.

This prototype, which the engineers nicknamed Betsy, was remarkably powerful but suffered from physical instabilities. The multiplier section would occasionally freeze mid-operation, trapping card decks inside. To clear the jam, the operators often resorted to physical force. In one famous incident, an engineer was told to "Kick it, Gib", and his literal kick drove a heavy metal cover directly into a 60-ampere fuse block, causing a massive, hazardous shower of sparks. When Northrop's founder, Jack Northrop, discussed the machine with IBM's CEO Thomas Watson to demand manufacturing support and standard parts for their modified system, IBM realized the commercial potential of the hybrid concept. IBM immediately flew Woodbury and his colleague George Fenn to New York to present their 603 and 405 combo. IBM's engineers then standardized the physical interfaces, replaced the 603 with the newer 604 calculating punch, and officially announced the commercial Card-Programmed Electronic Calculator in May 1949.

Even after the commercial release, Northrop continued to drive the CPC's evolution. Early users struggled because instruction cards had to point to highly specific, hardwired microprograms on a physical plugboard, making it almost impossible to share programs. To break this hardware bottleneck, Northrop's Rex Rice engineered a general-purpose control panel. By wiring a generalized set of logical pathways and math routing systems directly into the board, programmers could write various mathematical applications entirely on standard card decks without needing to manually rewire the plugboard for every new problem.

The ancestry of the CPC can even be traced a few years further back, to 1943 Los Alamos and the race to build the atomic bomb. The concepts and pieces that would become the CPC three years later were already being brought together at Los Alamos and even earlier at Columbia University, and it’s likely that some of the Northrop researchers had been present at Los Alamos and possibly Columbia. Here is a description of the arrival of the pieces of what soon would become known as the CPC at Los Alamos, and Richard Feynman’s reaction.

\begin{quotation}
Feynman, frustrated, turned to Nicholas Metropolis, a mustached Greek mathematician who later became an authority on computation and numerical methods, and said, “Let’s learn about these damned things and not have to send them to Burbank.” (Feynman grew a temporary mustache, too.) They spent hours taking apart new and old machines for comparative diagnosis; learned where the jams and slippages began; and hung out a shingle advertising, “Computers Repaired.” Bethe was not amused at this waste of his theoreticians’ time. He finally ordered a halt to the tinkering. Feynman complied, knowing that within weeks the shortage of machines would change Bethe’s mind. Escalation of the computation effort came in the fall of 1943 with an order to IBM for business machines to be delivered to an unknown location: three 601 multipliers, one 402 tabulator, one reproducer-summary punch, one verifier, one keypunch, one sorter, and one collator. Astronomers at Columbia had been experimenting with punch-card computing before the war. A multiplier, an appliance the size of a restaurant stove, could process calculations in large batches. Electrical probes found the holes in the cards, and operations could be configured by plugging groups of wires into a patchboard. Among the computation-minded at Los Alamos, the prospect of such machines caused excitement. Even before they arrived, one of the theorists, Stanley Frankel, set about devising improvements: for example, tripling the output by rearranging the plugs so that three sets of three- or four-digit numbers could be multiplied in a single pass. Having requisitioned the machines, the scientists now also requisitioned a maintenance man—an IBM employee who had been drafted into the army. They were gaining adroitness at military procurement. The crates arrived two days before the repairman; in those two days Feynman and his colleagues managed to get the machines unpacked and assembled, after a fashion, with the help of nothing but a set of wiring blueprints. \cite{levyHackersHeroesComputer1984}
\end{quotation}

This environment compelled Dantzig and Orchard-Hays to engage in a radical reinvention of Simplex, optimizing the algorithm to minimize the physical movement of data and the frequency of card-handling. The limitations of the CPC necessitated a rigorous focus on sparsity and numerical stability, as every additional iteration of the algorithm translated into significant manual labor and increased risk of mechanical error. By successfully executing complex linear programming problems on such primitive hardware, the RAND researchers established the pragmatic foundation for the field of numerical analysis. Their work proved that the efficacy of a model was as much a function of hardware-specific hacking as it was of theoretical elegance, setting a precedent for the nascent high-performance computing culture.
While Dantzig provided the mathematical vision, Orchard-Hays provided the computational genius. In 1952, when Dantzig tasked Orchard-Hays with implementing Simplex on the CPC, they quickly realized that the standard tableau form of Simplex, which updates the entire matrix at every step, was impossible on the CPC. When they first attempted an implementation, they did actually try to compute an explicit inverse at each iteration. Because the CPC relied on external punched cards rather than internal memory to hold the matrix, the machine operators were forced to physically re-run entire decks of cards through the accounting machine over and over. Dantzig stated that he was "appalled" by the results of this first implementation. 

This intellectual crisis forced them to propose the Revised Simplex Method in 1953. The machine simply did not have the memory to hold or update a large matrix of coefficients, and the numerical precision was nowhere near sufficient. Orchard-Hays transitioned the algorithm from a global matrix update to a more surgical, data-efficient procedure. This kept the original data static on one card deck and only computed the necessary column for the current step, maintaining the much smaller square matrix of the current basis inverse, either in the machine's limited storage or on a second smaller card deck. Dantzig recalled an idea that they named the Product Form of the Inverse. Instead of storing the inverse as a dense, explicit matrix, the PFI maintained the inverse as a product of elementary matrices known as Eta vectors. Eta vectors could be punched onto cards in the order they were generated. To update the current solution, the CPC simply had to read through the deck of Eta cards sequentially. Because most real-world linear programs are sparse, the PFI allowed them to write code that only stored and processed non-zero coefficients. This drastically reduced the burden on the CPC by minimizing the number of intermediate cards that operators had to punch and physically shuffle. By maintaining only the inverse of the current basis matrix as Eta vectors, while treating the original problem data as a static reference, they effectively bypassed the memory constraints of the CPC.
 
Recognizing that standard single-precision arithmetic would succumb to catastrophic rounding errors during the iterations of Simplex, Orchard-Hays engineered a software-hardware hybrid capable of something effectively like double-precision floating-point calculations. The IBM 604 performed arithmetic in Binary Coded Decimal. While it was possible to wire a CPC control panel for floating decimal operations, to handle the wide range of values encountered in linear programming the machine utilized a column shift unit that could physically move digits up to five positions. The 604 had to be painstakingly programmed via physical plugboards to handle the expanded word lengths. The resulting setup provided the numerical stability required, but it added the laborious work of changing the plugboard configurations. 
By decoupling the original dataset and exploiting sparsity, this second CPC implementation successfully solved linear programs with up to 45 constraints and 70 variables, including a complex model of the famous Stigler diet problem. The architectural shift allowed the CPC to solve significantly larger systems than its bare specifications would suggest, fundamentally altering the trajectory of numerical analysis. The Revised Simplex Method or Inverse Basis Method became a role-model for the concept of computational economy and the minimization of data movement and storage, principles that would become foundational for the field of high-performance computing. It was excruciatingly slow by modern standards. Solving a system with even twenty or thirty constraints could take an entire afternoon, requiring constant human supervision to move card decks between hoppers. Unlike the automated black box nature of modern software, running a program required the researchers to maintain a constant awareness of the machine's state, sensing the rhythmic cadences of the card reader that signaled a successful iteration, and improvements took concrete physical form simply by reducing the physical burden of card-shuffling. 

The work done at RAND in 1952 was epochal for numerical modeling and optimization. It moved linear programming from a manual process, done by human computers with desk calculators, to an automated one. Before 1952, some of the earliest tests of Simplex on a 77-variable problem had required approximately 120 man-days using hand-operated desk calculators. By transitioning to an automated sequence on the IBM CPC, Dantzig and Orchard-Hays demonstrated that complex models could be handled purely by machine logic. This period established the critical precedent that models and algorithms must be meticulously tailored to the specific architectures of their host hardware. This hardware-first approach to software design ensured that high-level strategic models could be executed within the stark physical limitations of the early digital age. The success of these early automated trials provided the necessary empirical validation for the first generation of commercial linear programming codes.

Driven by the massive economic incentives of optimizing refinery yields and gasoline blending, the oil industry quickly became a primary adopter and financial subsidizer of this new technology. By the sixties, the software lineage started by Orchard-Hays evolved into codes like LP/90, which ran on the IBM 7090 and could successfully solve models containing over 1,000 constraints. Today’s LP models regularly have millions of variables and constraints. As industry moved from the electro-mechanical CPC to the fully electronic IBM 701 and the subsequent 704, the first machine with built-in floating-point hardware, the methodologies forged at RAND were rapidly ingrained in the modeling culture. These true computers allowed for the scaling of Dantzig's work, enabling the oil and aerospace industries to move beyond small-scale problems toward the massive systems of equations that define modern industrial life. 

As will be discussed in greater detail later, UT Austin inherited an IBM CPC from Humble Oil in Houston, later Exxon, in 1958. This was made possible by Frederick Matsen, a pioneering professor of chemistry and physics. The installation of the machine was a remarkably grassroots effort. According to former UT President and fellow chemistry professor Norman Hackerman, interested faculty members and graduate students literally carried the IBM CPC into Welch Hall and installed it themselves. They chose to operate the machine with minimal involvement from higher authorities specifically to avoid excessive paperwork. While Matsen had already been a pivotal force in pushing UT administration to embrace computing as early as 1950, the arrival of the IBM CPC accelerated computational research on campus. The machine's presence clearly demonstrated the value of computational power, leading to a continuous demand for greater capacity. Matsen's efforts played an important role in the formal establishment of the UT Computation Center and the eventual acquisition of the CDC 6600 supercomputer.

\section{Exxon And Bill Drews}

The collaboration between George Dantzig and Exxon’s Bill Drews in the sixties helped strengthen the ongoing technology transfer of computational modeling from the defense sector to global industry. While RAND had proven the technical viability of Simplex, Drews possessed a rare ability to bridge the gap between Dantzig’s models and the physical realities of industry. The inherent nature of refinery operations, characterized by continuous liquid flows and complex chemical constraints, offered a natural domain for the application of linear programming. In 1962 Drews took a leave of absence to study under Dantzig at UC Berkeley and soon afterward Drews and Leo Rapoport published a landmark paper in the Harvard Business Review titled "Mathematical Approach to Long-Range Planning" that shifted the focus of LP toward 10-to-30-year strategic capital investments and financial modeling. Drews himself stated that Exxon used linear optimization to "schedule our tanker fleets, design port facilities, blend gasoline, create financial models, you name it". This holistic approach treated the massive organization as a mathematically optimized system.

As early as 1952, Simplex was already in use by the oil industry. Gulf Oil is credited with the earliest formally recognized application of LP to the blending problem. In 1949, Gulf researchers William Cooper and Bob Mellon were struggling with the problem of octane giveaway, the expensive practice of using excessive high-quality additives to ensure aviation fuel met minimum standards. Following a serendipitous meeting with mathematician Abraham Charnes, who applied Dantzig’s newly published Simplex method, the team successfully developed an LP formulation in 1952 that balanced the mixture's ignition, volatility, and octane requirements at the lowest possible cost.

The transition from recipe-based petroleum refining to mathematically optimized processing marked a paradigm shift in the economic structure of the Texas oil industry. Historically, refineries operated on fixed ratios, pre-determined recipes that ensured product quality but failed to account for the marginal costs of various crude inputs. Prior to linear programming, refinery optimization was conducted through tedious hand balances and trial-and-error to find acceptable operating conditions. Linear programming replaced this guesswork, allowing engineers to dynamically select the best combination of feedstocks and processing methods to maximize the gross refining margin, the difference between product revenues and operational costs. 

The 1955 publication of Linear Programming: The Solution of Refinery Problems by G.H. Symonds signaled that Exxon had moved beyond experimental trials toward a model of data-driven industrial governance. This era effectively institutionalized Operations Research within the energy sector, providing a competitive advantage that fundamentally altered the economic landscape in Texas. Exxon’s pioneering adoption of Operations Research techniques extended far beyond simple process optimization, encompassing the entire lifecycle of the petroleum supply chain. 

By applying Dantzig’s algorithms to the management of oil field output, it was possible to maximize short-term yield without violating the physical constraints of the reservoir, preventing long-term damage to the field’s pressure and productivity. The use of LP for vessel scheduling and port design transformed the logistics of oil distribution into a precise science. This holistic approach to systems engineering established a technocratic precedent. The refinery was no longer viewed as a collection of disparate machines, but as a unified, mathematically optimized system, and this Exxon Model served as the foundational blueprint for the industrial computing programs that led to the state's role as a global leader in both energy and computational science.

In the late fifties and early sixties, Exxon pioneered the application of LP to the blending problem, allowing for the dynamic adjustment of refinery inputs based on real-time market data. This computational approach enabled engineers to maintain rigorous specifications for octane and viscosity while simultaneously minimizing the cost of expensive chemical additives. By the seventies, this optimization framework was extended to the complex logistics of daily refinery operations, largely through the modeling innovations of researchers like T.E. Baker at Exxon. As global energy markets became increasingly volatile, the ability to schedule refinery throughput with mathematical precision became a critical competitive advantage. Baker moved beyond basic blending to manage the flow of materials over time. He later built upon this foundation to develop MIMI, Manager for Interactive Modeling Interfaces, a highly configurable system that integrated LP with database management and graphical scheduling functions. Systems like MIMI allowed refineries to synchronize the arrival of crude with distillation schedules and product distribution networks. These tools established a new standard for operational scheduling, a methodology that was rapidly adopted across the global energy sector. This era of high-performance industrial modeling provided the intellectual and financial capital that fueled the growth of specialized numerical analysis research programs. 

Bill Drews’ evolution from a chemical engineer to the intellectual leader of Exxon’s Operations Research group represents the moment when LP transcended its tactical origins to become a foundational instrument of corporate governance. He started as a chemical engineer in Research and Engineering and held patents in areas such as catalyst pretreatment. His subsequent study under George Dantzig allowed him to reconceptualize the entire corporate entity as a unified modeling system. By the early sixties, Drews had shifted his focus toward long-range planning, developing models that projected Exxon’s operational and financial trajectories decades into the future. His dual role as a practitioner at Exxon and a referee for the journal Operations Research gave him a unique authority to define the technical standards of the field, ensuring that LP would serve as the primary blueprint for the emergence of modern systems engineering. 

The publication of "Mathematical Approach to Long-Range Planning" in 1962 was a landmark in this new technocratic paradigm. It definitively broadened the applications of LP from short-term tactical scheduling, like aviation gasoline blending, toward comprehensive, long-term strategic investments. The new models allowed executives to plan operations over a multi-decade horizon. Instead of optimizing a single chemical process, they evaluated a massive set of interconnected variables to minimize total discounted costs while satisfying intertemporal capacity balances, environmental protection standards, and overall system reliability constraints. The paper demonstrated that LP was a framework for managing the vast complexities of global investment and resource procurement. By modeling a corporation’s future as a set of interconnected linear inequalities, Drews and Rapoport provided executives with a mechanism to quantify the long-term consequences of current decisions. This methodology eventually proliferated beyond the oil industry, most notably into aerospace, automotive, and the airlines. 

The recurring consultancy of Dantzig at Exxon underscored the corporation’s role as a modeling leader during the fifties and sixties, while many companies were still performing calculations by hand. Confronted with logistical challenges that seemed unsolvable using contemporary computational methods, Exxon’s Operations Research group leveraged Dantzig’s expertise to develop sophisticated decomposition techniques. The Decomposition Principle was developed by Dantzig and Philip Wolfe in 1959 and 1960. This breakthrough technique allowed extremely large linear programs with special structures to be broken down and solved in parts, which was essential for modeling the massive organizational systems and logistical architectures of the era. Drews frequently hired Dantzig as a consultant, and he would fly in to help the Exxon team break down global problems and bypass the memory and processing limitations of early mainframes by partitioning global models into manageable sub-systems. By funding this level of high-theory research, Exxon effectively bridged the divide between the LP models formulated at RAND and the pragmatic, high-stakes requirements of industry. The Dantzig-Drews collaboration also bridged the gap between corporate planning and pure macroeconomic theory. Drews communicated an insight to Dantzig that in a competitive economy, the distribution of money flows is often arbitrarily dictated by political processes rather than the invisible hand of pure market forces. This led Dantzig to formulate and mathematically prove the Drews Institutionalized Divvy Economy model in 1973. 

The development of internal economic simulations at Exxon during the sixties serves as a definitive precursor to modern decentralized resource management. By assigning shadow prices to internal assets, Dantzig and Drews enabled Exxon to operate as a mathematically coordinated marketplace rather than a rigid hierarchy. These shadow prices, derived from the dual values of LP models, provided a rigorous mechanism for capital allocation, forcing competing divisions, such as refining and supply, to purchase resources at prices that reflected their true marginal value to the total system. It provided a mathematical mechanism for corporate headquarters to coordinate the activities of decentralized divisions by passing down prices rather than direct operational commands. In this model, corporate headquarters acts as a coordinator. To ensure that independent divisions act in the best interest of the entire firm, the coordinator uses optimal shadow prices to essentially charge the divisions for their use of shared corporate resources. This approach allowed executive leadership to maintain strategic coherence without engaging in the granular micromanagement of local operations. 
Beyond the immediate benefits of internal market simulation, the technical challenges Exxon faced in managing its global models precipitated a revolution in numerical analysis, specifically in the field of sparse matrix computation. Because a global supply chain model consists primarily of localized interactions, the resulting mathematical matrices are sparse, containing a high percentage of zero-valued entries. Early LP models heavily relied on exploiting the fact that most large-scale matrices contain a high percentage of zero-valued entries. The foundational algorithmic techniques for this were pioneered at RAND. Orchard-Hays implemented the first commercial-grade linear programming software at RAND in 1954, incorporating the first theoretical ideas on how to compact the inverse and take advantage of sparsity. Harry Markowitz, another RAND luminary, introduced a pivot selection rule in 1957 specifically designed for inverting matrices in LP. His technique became widely adopted for solving massive systems of simultaneous equations where the coefficients are mostly zero. 

While Exxon's researchers did not invent sparse matrix computation, they did play a vital financial role in scaling it. In the early 1970s, as models grew increasingly massive, developer James Kalan introduced the concept of supersparsity. Exxon directly supported Ketron Corporation in using Kalan's work to develop the MPS III software system. Tuned for the new IBM 360 platforms, this software represented a quantum leap in speed and problem size, utilizing advanced sparsity handling to solve models with up to 32,000 constraints.

To further this work, the company formally established Mathematics and Systems in 1967, a unit largely devoted to advancing operations research and information technology. At the time, commercial software providers were primarily focused on record-keeping and basic ledger functions. They were entirely incapable of processing the massive datasets and non-linear chemistry required to optimize the global trajectory of hydrocarbons from Middle Eastern extraction sites, through European refineries, to North American retail networks. With the formalization of its OR efforts in the Mathematics and Systems group in the late sixties, Exxon established an intellectual ecosystem that rivaled the top academic departments in the country. The corporation became a second home for the field’s elite practitioners, who utilized Exxon’s massive industrial footprint as a real-world laboratory for testing the limits of algorithmic stability. M\&S hosted a diverse array of luminaries, from computing pioneers like Grace Murray Hopper and Bob Bemer to control theory experts like Dimitri Bertsekas. 

The proprietary matrix-handling and large-scale system decomposition techniques they pioneered rapidly permeated academic circles, effectively setting the standard curriculum for engineering departments nationwide. The methodologies developed through the Dantzig-Exxon partnership, particularly in the realm of large-scale system decomposition, bridged the gap between academic theory and real-world application. The decomposition principle allowed the company to bypass mainframe hardware limits by breaking its massive global supply chain into independent subproblems, such as individual refinery operations, coordinated by a master problem that handled overarching linking constraints like pipeline capacity and tanker availability. Furthermore, to test the limits of algorithmic stability in real-time production optimization, M\&S practitioners utilized piecewise linearization. This technique converted highly non-linear physical challenges, such as gas-coning in horizontal wells on the Norwegian Continental shelf, into scalable mixed-integer linear programming models that could be solved rapidly. 
All of this was equally true for the Exxon Production Research group and the Buffalo Speedway Campus in Houston, a legendary site in the history of privately funded research. While the Mathematics and Systems group focused on LP and OR, the Production Research group focused on subsurface modeling and simulation. Production Research is where the company housed some of the most powerful privately owned supercomputers in the world and transformed the oil industry from a field of empirical exploration into a rigorous discipline of digital simulation, establishing Houston as the global epicenter of computational geophysics. 

For decades, the heart of this work was located at Buffalo Speedway, a research facility where geophysicists and computer scientists worked in a university-like atmosphere similar to Bell Labs and, at the time, possibly the most important privately-owned earth science research center in the world. This is where three-dimensional seismic imaging, reservoir simulation, and deep-water drilling technologies were invented. In the sixties, researchers pioneered the concept of three-dimensional seismic technology. Before this, geologists focused on two-dimensional slices of the earth. Because early mainframes lacked the memory to hold entire 3D seismic volumes, Exxon researchers had to invent clever data-shuffling techniques and proprietary codes to process information in smaller chunks without losing the spatial context of the larger geological structure. 

All of this created a critical market for the high-performance computing industry, driving companies like Cray to build faster processors and interconnects specifically for Exxon. Exxon’s researchers pioneered the algorithmic techniques necessary to integrate disparate slices into a coherent 3D data cube, a process that required unprecedented levels of data throughput and memory management. This revolution in imaging did more than simply improve drilling success rates. It established the foundational requirements for the next generation of supercomputing architecture. The necessity of processing these massive seismic volumes provided the primary industrial justification for the development of vector processing and parallel computing.

Beyond seismic imaging, Buffalo Speedway researchers also established the digital simulation foundations for fluid extraction. By solving complex fluid flow equations across millions of grid cells, they created digital twins of oil fields to optimize the placement of wells and the injection of water or gas. As the company pushed into offshore environments exceeding 5,000 feet in depth, these simulations guided the development of the specialized materials, riser systems, and blowout preventers necessary to operate in high-pressure, low-temperature environments. The Buffalo Speedway campus functioned as a terrestrial analog to the premier aerospace laboratories of the Cold War, united by a shared reliance on the finite-difference method to model complex physical phenomena. Geophysicists utilized FDM to achieve the spatial and temporal discretization of the partial differential equations that govern fluid flow and seismic wave propagation. This process effectively translated the continuous physical properties of the earth’s crust into massive, sparse systems of linear equations. 

By partitioning a reservoir into a digital grid, engineers could simulate pressure transients and saturation fronts with high fidelity, a task that required the deployment of advanced preconditioned iterative solvers. This rigorous computational approach transformed reservoir engineering from a descriptive craft into a predictive science, mirroring the methodological evolution of computational fluid dynamics in the aerospace industry. The resulting triangle of research, connecting Buffalo Speedway with Rice University and UT Austin, established a regional hegemony in the field of high-performance computing. This relationship was fueled by geographic proximity and a mutual obsession with the inverse problem of seismic imaging. Looking ahead, it’s useful to note that seismic imaging stands beside gravity modeling and weather forecasting as foundational members of the class of historic inverse problems, where observational effects are inverted to model hidden causes, whether within the Earth or the Earth’s atmosphere. This is a recurring theme for computational modeling and simulation.

The Buffalo Speedway campus was a few miles from Rice, allowing for a unique, porous relationship, and Rice has historically been a leader in computing as a result. At the same time, the pipeline between the campus and UT was massive. UT is frequently one of the top feeder schools for Exxon, and at one point the company employed over 2,500 UT alumni. The alumni network within the Buffalo Speedway campus was so strong that it influenced the company's internal culture. UT naturally developed its strengths in computational science around FDM, as will be shown in the following sections tracing the stories of David Young and the establishment of the UT Computation Center. Exxon’s strategic funding of fellowships and research at Rice and UT allowed for a unique exchange of intellectual talent, where academic breakthroughs in numerical analysis were immediately pressure-tested against Exxon’s massive seismic datasets. This collaborative ecosystem played a major role in the formation of the algorithmic bedrock of modeling and simulation. 
